\documentclass[11pt]{article}

\usepackage[a4paper,margin=1in]{geometry}
\usepackage{fontspec}
\setmainfont{Calibri}
\usepackage[hidelinks]{hyperref}
\usepackage{enumitem}
\usepackage{listings}
\usepackage{xcolor}

\lstset{
  basicstyle=\ttfamily\small,
  breaklines=true,
  columns=fullflexible,
  frame=single,
  showstringspaces=false
}

\title{dstr: A Small Language for Systems That Become Graphs}
\author{Sambuddha Majumder \and Jayanta Majumder}
\date{\today}

\begin{document}
\maketitle

\begin{abstract}
\texttt{dstr} is a compact language for describing finite dynamic systems in a
form that is both easy to write and semantically explicit. Its primary
interface is a small s-expression DSL in which one states variables, initial
conditions, actions, invariants, and reachability goals without the verbosity
that often discourages exploratory modeling. The resulting description is not
intended merely to support a checker run. It is intended to generate an
explicit state graph that can subsequently be queried, transformed, filtered,
and published. The central claim of this paper is that a state-action language
becomes substantially more valuable when its semantics are treated as a
first-class graph artifact, suitable not only for validation but also for
downstream graph analysis and presentation.
\end{abstract}

\section{Introduction}

The name \texttt{dstr} was chosen partly for its sound. It echoes the name
Dijkstra, and in fact retains the non-silent consonants from that name. For
what it is worth as an expansion, it also fits the phrase ``describing states
and transitions for reasoning''. The name is therefore meant to suggest both a
style of formal thought and the specific structural emphasis of the language.

Many important systems are easier to understand as spaces of possibilities than
as streams of execution. A protocol, a controller, a workflow, or a concurrent
routine is not merely a program that runs. It is a set of states, a family of
allowed transitions, and a collection of claims we hope remain true throughout
motion. Once viewed in that way, the natural mathematical object is a graph.

Yet graph construction is usually treated as a downstream detail. One first
writes code, or a large formal model, or an ad hoc simulation, and only later
attempts to recover a graph worth studying. The result is often unsatisfactory.
Either the notation is pleasant but semantically loose, or the semantics are
precise but the authoring burden is too high for exploratory use.

\texttt{dstr} begins from a different preference. The system author should be
able to write down state and action structure in a small, expressive, tree-like
surface language. That description should compile into a simple internal form,
be checked by explicit exploration, and then be materialized as a graph that is
valuable in its own right. The graph is not merely evidence that checking took
place. It is a reusable object for later analysis.

\section{The High-Level Interface}

The primary interface of \texttt{dstr} is a Lisp-style DSL. This choice is not
an affectation. S-expressions provide exactly the right compromise for a system
description language:

\begin{itemize}[leftmargin=*]
\item they are compact,
\item they expose tree structure directly,
\item they admit macros and normalization naturally, and
\item they avoid the visual clutter of deeply tagged abstract syntax.
\end{itemize}

The aim is that one writes the system in a form close to its conceptual shape.
State variables are named directly. Predicates appear as ordinary nested forms.
Action structure is explicit without becoming ceremonial. The next-state view of
a variable can be expressed tersely, and symbols can serve as literals where
that is the most natural reading.

The DSL is also intentionally programmable. A specification file may define
local macros and helper functions, or load a small library of such definitions,
before stating the system itself. Ordinary expansion can be requested
explicitly, and the current front end also supports compact macro shorthands in
the surface notation. The implemented surface forms are precise. One may write
\texttt{(expand macro-name ...)} for an explicit one-step expansion,
\texttt{(\%same x)} as shorthand for \texttt{(expand same x)}, and
\texttt{(!unchanged x y)} when the macro expands to a list of sibling clauses
to be spliced into the surrounding form. The predefined helper
\texttt{same} expands to \texttt{(= x+ x)}, while \texttt{unchanged} expands to
a list such as \texttt{((= x+ x) (= y+ y))}. The front end also accepts
top-level \texttt{defmacro}, \texttt{defun}, and \texttt{load} forms before the
\texttt{system} form, so a specification can import a macro library and use
ordinary Lisp code at expansion time. Such macros are useful as genuine
shorthands rather than mere decoration: they allow families of related actions,
invariants, or structural patterns to be written once and reused declaratively
across a model.

Two small conventions carry much of the language. First, an unadorned variable
name such as \texttt{light} or \texttt{lock} denotes that variable's value in
the current state. A name with the suffix \texttt{+}, such as \texttt{light+}
or \texttt{lock+}, denotes its value in a successor state. Second, an
\texttt{action} clause is not a command body in the ordinary programming sense.
It is a predicate relating current-state values and next-state values. One does
not instruct the machine to execute steps imperatively. One characterizes which
pairs of states count as valid instances of the named transition. In the
\texttt{next} clause, a form such as \texttt{@turn-on} refers to the named
action predicate \texttt{turn-on}; larger \texttt{next} formulas can combine
such references to define the overall transition relation.

This matters more than syntax aesthetics. In a language for behavioral models,
verbosity is not a cosmetic flaw. It actively suppresses exploration. If every
small change requires plumbing through a large ceremonial representation, users
stop trying variants, stop expressing auxiliary properties, and stop thinking in
terms of the full state space. A compact DSL preserves the willingness to model.

In \texttt{dstr}, the high-level interface is not a veneer over an opaque core.
Surface forms are normalized into a clear intermediate representation and then
emitted as JSON when needed. This preserves a concise authoring language while
retaining a stable machine-oriented form for checking, tooling, and regression
tests. The separation is important: it allows the front end to remain pleasant
without sacrificing semantic clarity or downstream interoperability.

The macro layer strengthens this separation rather than weakening it. Authors
are free to use abstraction aggressively at the surface level, but those
abstractions are compiled away before checking. The checker therefore continues
to operate on an explicit normalized representation even when the source model
was assembled from reusable local idioms. This makes the language feel much
smaller to write than the eventual transition relation it denotes.

\subsection{Macro Use in Practice}

A representative use of the macro system is to load a small library of helper
macros, define expansion-time helper functions, and then generate whole
families of action clauses rather than writing out near-duplicates by hand. The
\texttt{dstr} distribution includes examples of this style in its test suite.
The following excerpt is representative:

\begin{lstlisting}[language=Lisp]
(load "cas-macro-lib.lisp")

(defun initial-to-try (pc)
  `(and
    (= ,pc a0)
    (= ,(next-var-symbol pc) 'try)))

(defmacro start-attempt (pc other-pc)
  (cons (initial-to-try pc)
        (preserve-clauses other-pc 'mem)))

(defmacro p1-actions ()
  '((action p1-start
      (!start-attempt p1pc p2pc))
    (action p1-cas-success
      (= p1pc try)
      (= mem 0)
      (= p1pc+ won)
      (!preserve p2pc)
      (= mem+ 1))))
\end{lstlisting}

Several distinct macro mechanisms appear here. The helper macro
\texttt{preserve} expands into a list of frame-condition clauses. The form
\texttt{(!preserve p2pc)} therefore splices \texttt{(= p2pc+ p2pc)} directly
into an action body. The local macro \texttt{start-attempt} combines a helper
function with library code to produce a reusable transition pattern, and
\texttt{(!p1-actions)} or \texttt{(!p2-actions)} can then splice whole groups
of generated \texttt{action} clauses into the surrounding \texttt{system}. This
is not cosmetic syntax sugar. It is a way to factor out repeated state-update
structure while still compiling to the same explicit normalized transition
predicates used by the checker.

\section{A Language for State and Action}

The conceptual center of \texttt{dstr} is straightforward: a system is described by
variables ranging over finite domains, a predicate selecting the initial
states, a set of named actions, and a transition relation assembled from those
actions. Around that core sit invariants and reachability-oriented properties.

This point is worth stating explicitly. In \texttt{dstr}, an action is written
as a predicate over two states: the current state and a candidate successor
state. Conditions involving plain variable names constrain the source state;
conditions involving the corresponding \texttt{+}-suffixed names constrain the
target state. The \texttt{next} clause then specifies which action predicates,
or which combinations of them, are admitted as actual transitions of the
system. Model checking proceeds by finding state pairs that satisfy those
predicates, starting from states that satisfy \texttt{init}, and assembling the
reachable transition graph from there. This organization follows the familiar
state-machine style exemplified in TLA+, where one separates an initial-state
predicate from next-state actions over current and successor variables
\cite{DBLP:books/aw/Lamport2002}.

Equally important is the order in which the semantics should be understood.
First, the state universe is determined as the Cartesian product of the
declared variable domains. Second, each action denotes a predicate over a pair
of states: a current state and a candidate successor state. Third, the
\texttt{next} clause combines these action predicates into a single transition
relation on the state universe. Finally, reachability from \texttt{init}
determines the actual state-transition graph of the system. This way of stating
the semantics is clarifying because it separates the space of possible states,
the logical definition of allowed transitions, and the reachable behavior that
is ultimately explored.

This style of description has an important philosophical consequence. It asks
the author to define behavior intensionally rather than procedurally. One says
what constitutes a legal state and what changes are allowed, not which control
path a single run happened to follow. That shift in viewpoint is exactly what
makes enumeration meaningful. Once a system is expressed as a finite state and
action theory, the reachable space can be explored completely.

Explicit enumeration is sometimes dismissed as unsophisticated because it does
not aspire to hide the state space behind symbolic compression. For the class of
systems targeted by \texttt{dstr}, that criticism misses the point. In early
design work and in many finite operational models, explicitness is an advantage.
It gives exact reachable states, exact transitions, exact deadlocks, and exact
counterexample paths. Most importantly, it yields a graph that can be inspected
and repurposed without semantic ambiguity.

\section{Examples}

Examples clarify the intended use of \texttt{dstr} more effectively than
abstract description alone. The following cases are not included merely as
syntax samples. They illustrate the form of feedback produced by the model
checker and the way in which that feedback supports design understanding. In
each example, occurrences of a variable such as \texttt{x} refer to the current
state, while \texttt{x+} refers to the successor state being constrained by the
action predicate.

\subsection{A Minimal Switching System}

The following specification describes a two-state switch:

\begin{lstlisting}[language=Lisp]
(system light-switch
  (vars light)
  (domain light off on)
  (init (= light off))
  (action turn-on
    (= light off)
    (= light+ on))
  (action turn-off
    (= light on)
    (= light+ off))
  (next (or @turn-on @turn-off))
  (invariant type-ok
    (in light (set off on)))
  (property eventually-on
    (eventually (= light on))))
\end{lstlisting}

This is a tiny system, but it demonstrates the basic rhythm of the language.
There is one variable, two actions, a simple type discipline expressed as an
invariant, and a reachability property. When model checked, the current
implementation reports:

\begin{quote}
Spec: light-switch\\
Reachable states: 2\\
Transitions explored: 2\\
Properties: \{eventually-on=true\}
\end{quote}

Even here the report is useful. It confirms that the intended state space is
exactly the one the author had in mind, that the transition relation is neither
over-constrained nor under-constrained, and that the desired state is indeed
reachable. The resulting graph is small enough to be grasped as a whole, which
makes it a useful introductory example both for teaching and for validating the
surface notation itself.

\subsection{A Resource Protocol}

A slightly richer example models two processes competing for a lock. In
\texttt{dstr}, one writes the actions that acquire and release the lock and then
asks for both safety and reachability:

\begin{lstlisting}[language=Lisp]
(system mutex-2proc
  (vars p1 p2 lock)
  (domain p1 idle cs)
  (domain p2 idle cs)
  (domain lock free p1 p2)
  (init (= p1 idle) (= p2 idle) (= lock free))
  (action p1-enter
    (= p1 idle) (= lock free)
    (= p2 p2+)
    (= p1+ cs) (= lock+ p1))
  (action p1-exit
    (= p1 cs) (= lock p1)
    (= p2 p2+)
    (= p1+ idle) (= lock+ free))
  (action p2-enter
    (= p2 idle) (= lock free)
    (= p1 p1+)
    (= p2+ cs) (= lock+ p2))
  (action p2-exit
    (= p2 cs) (= lock p2)
    (= p1 p1+)
    (= p2+ idle) (= lock+ free))
  (next (or @p1-enter @p1-exit @p2-enter @p2-exit))
  (invariant mutual-exclusion
    (not (and (= p1 cs) (= p2 cs))))
  (property p1-can-reach-cs (eventually (= p1 cs)))
  (property p2-can-reach-cs (eventually (= p2 cs))))
\end{lstlisting}

The checker currently reports:

\begin{quote}
Spec: mutex-2proc\\
Reachable states: 3\\
Transitions explored: 4\\
Properties: \{p1-can-reach-cs=true, p2-can-reach-cs=true\}
\end{quote}

This is the sort of feedback that matters in design. The invariant states that
the two processes are never simultaneously in the critical section. The
properties establish that each process can still reach that section. Together
they distinguish mere exclusion from useful exclusion. The graph then makes the
structure behind the report visible: an idle state and two ownership states,
with the lock mediating the movement between them. Even in this small example,
one can see how a behavioral design becomes a structured object rather than a
prose claim.

\subsection{A Deliberately Broken Design}

The value of a state-action language is equally clear when the design is wrong.
Consider a pathological ``mutual exclusion'' protocol in which one action
places both processes into the critical section at once:

\begin{lstlisting}[language=Lisp]
(system broken-mutex
  (vars p1 p2)
  (domain p1 idle cs)
  (domain p2 idle cs)
  (init (= p1 idle) (= p2 idle))
  (action both-enter
    (= p1+ cs)
    (= p2+ cs))
  (next @both-enter)
  (invariant mutual-exclusion
    (not (and (= p1 cs) (= p2 cs)))))
\end{lstlisting}

Here the checker does not simply say that the invariant fails. It shows how:

\begin{quote}
Spec: broken-mutex\\
Reachable states: 2\\
Transitions explored: 2\\
Properties: \{\}\\
Invariant violations:\\
\ \ - mutual-exclusion path=[\{p1=idle, p2=idle\}, \{p1=cs, p2=cs\}]
\end{quote}

That path is exactly the sort of evidence a designer needs. It is short,
concrete, and already graph-shaped. One can mark it in a visualization, isolate
the error boundary with graph tools, or compare it to a repaired design. This
illustrates why explicit-state output is valuable beyond verdicts. It supports
diagnosis.

\subsection{A Small Search Problem}

\texttt{dstr} is not limited to protocol-like examples. It also handles finite
search spaces naturally. The following model captures a water-jug puzzle with a
five-gallon jug and a three-gallon jug:

\begin{lstlisting}[language=Lisp]
(system water-jugs
  (vars big small)
  (domain big 0 1 2 3 4 5)
  (domain small 0 1 2 3)
  (init (= big 0) (= small 0))
  (action fill-small-jug
    (= small+ 3)
    (= big+ big))
  (action fill-big-jug
    (= big+ 5)
    (= small+ small))
  ...
  (property can-measure-four-gallons
    (eventually (= big 4))))
\end{lstlisting}

The checked version in the repository includes the expected pouring and emptying
actions. The current model checker reports a reachable space of 16 states,
explores 74 transitions, and confirms that the four-gallon target is reachable:

\begin{quote}
Reachable states: 16\\
Transitions explored: 74\\
Properties: \{can-measure-four-gallons=true\}
\end{quote}

This example is useful because it is neither purely didactic nor obviously
protocol-shaped. It shows that the language can describe constrained
transformations over a finite state space and answer a concrete goal question:
is the target configuration reachable? Once the graph exists, one can move
beyond the yes-or-no answer and ask which subgraph contains all successful
strategies, which states are bottlenecks, and how many distinct routes lead to
the target.

\subsection{A Three-Process Bakery Protocol}

A more substantial example is a bounded three-process model of the bakery
algorithm. The purpose of including it is twofold. First, it shows that
\texttt{dstr} can express a recognizably algorithmic coordination protocol in a
 direct and readable style. Second, it shows what verification means in a case
where the model is large enough to be interesting, yet still explicit enough to
be fully explored.

The model uses, for each process, a control-state variable, a boolean
\texttt{choosing} flag, and a ticket number. In the bounded finite model used
here, ticket values range from \texttt{0} to \texttt{4}. The control states are
\texttt{a0} (idle), \texttt{a1} (choosing a ticket), \texttt{wait}, and
\texttt{cs}. The following excerpt shows the overall structure together with the
first process's actions:

\begin{lstlisting}[language=Lisp]
(system bakery-3proc
  (vars p1pc p2pc p3pc ch1 ch2 ch3 n1 n2 n3)
  (domain p1pc a0 a1 wait cs)
  (domain p2pc a0 a1 wait cs)
  (domain p3pc a0 a1 wait cs)
  (domain ch1 t nil)
  (domain ch2 t nil)
  (domain ch3 t nil)
  (domain n1 0 1 2 3 4)
  (domain n2 0 1 2 3 4)
  (domain n3 0 1 2 3 4)
  (init
    (= p1pc a0) (= p2pc a0) (= p3pc a0)
    (= ch1 nil) (= ch2 nil) (= ch3 nil)
    (= n1 0) (= n2 0) (= n3 0))

  (action p1-start
    (= p1pc a0)
    (= p1pc+ a1)
    (= ch1+ t)
    (= p2pc+ p2pc) (= p3pc+ p3pc)
    (= ch2+ ch2) (= ch3+ ch3)
    (= n1+ n1) (= n2+ n2) (= n3+ n3))

  (action p1-choose
    (= p1pc a1)
    (= p1pc+ wait)
    (= ch1+ nil)
    (or
      (and (>= n2 n3) (< n2 4) (= n1+ (+ n2 1)))
      (and (> n3 n2) (< n3 4) (= n1+ (+ n3 1))))
    (= p2pc+ p2pc) (= p3pc+ p3pc)
    (= ch2+ ch2) (= ch3+ ch3)
    (= n2+ n2) (= n3+ n3))

  (action p1-enter
    (= p1pc wait)
    (not ch2)
    (not ch3)
    (or (= n2 0) (< n1 n2) (= n1 n2))
    (or (= n3 0) (< n1 n3) (= n1 n3))
    (= p1pc+ cs)
    (= p2pc+ p2pc) (= p3pc+ p3pc)
    (= ch1+ ch1) (= ch2+ ch2) (= ch3+ ch3)
    (= n1+ n1) (= n2+ n2) (= n3+ n3))

  (action p1-exit
    (= p1pc cs)
    (= p1pc+ a0)
    (= n1+ 0)
    (= ch1+ nil)
    ...)

  (next
    (or @p1-start @p1-choose @p1-enter @p1-exit
        @p2-start @p2-choose @p2-enter @p2-exit
        @p3-start @p3-choose @p3-enter @p3-exit))

  (invariant mutual-exclusion
    (not
      (or (and (= p1pc cs) (= p2pc cs))
          (and (= p1pc cs) (= p3pc cs))
          (and (= p2pc cs) (= p3pc cs)))))

  (property p1-can-reach-cs (eventually (= p1pc cs)))
  (property p2-can-reach-cs (eventually (= p2pc cs)))
  (property p3-can-reach-cs (eventually (= p3pc cs))))
\end{lstlisting}

Several features of the language are visible here. The actions are not written
as assignment blocks but as predicates over current and successor states.
Unchanged variables are stated explicitly, which makes each action a complete
binary-state condition. The \texttt{p1-choose} action computes a bounded version
of ``one plus the maximum competing ticket'', while \texttt{p1-enter} encodes
the waiting condition that the other processes are not choosing and do not hold
a lexicographically prior ticket.

The checker explores this bounded model successfully. It reports:

\begin{quote}
Spec: bakery-3proc\\
Reachable states: 200\\
Transitions explored: 390\\
Properties: \{p1-can-reach-cs=true, p2-can-reach-cs=true,
p3-can-reach-cs=true\}
\end{quote}

Most importantly, the mutual-exclusion invariant is preserved throughout the
reachable state space. The model also includes two simple structural invariants:
that an idle process has ticket \texttt{0}, and that a process can have its
\texttt{choosing} flag raised only while in the ticket-selection control state.
These invariants likewise hold over all reachable states.

There is, however, an important modeling caveat. Because the ticket domain is
bounded artificially to keep the model finite, the checker also finds deadlocks
in states where a process would need a ticket value beyond the chosen bound.
This does not contradict the safety result. It reflects the fact that the model
is a finite approximation of the unbounded bakery protocol. In this setting,
\texttt{dstr} is verifying a meaningful safety claim under a deliberate state
space restriction, while simultaneously exposing where that restriction begins
to matter operationally.

\subsection{What the Reports Mean}

Across these examples, the current checker reports four kinds of information of
lasting value:

\begin{itemize}[leftmargin=*]
\item the number of reachable states, which tells us the effective behavioral
size of the system,
\item the number of explored transitions, which gives a first sense of the
system's dynamical richness,
\item property results, which answer the author's stated reachability
questions, and
\item invariant violations with concrete paths, which turn a failed design into
an inspectable artifact.
\end{itemize}

The importance of these reports lies not only in certifying or refuting claims.
They also prepare the model for its second life as a graph. Counts orient the
reader, property results identify semantically interesting regions, and
counterexample paths provide anchors for subsequent graph slicing and analysis.

\section{Why the Graph Matters}

The state graph produced from a \texttt{dstr} model is the central artifact of
the method. It records not only what is possible, but how possibility is
organized. Which regions are tightly connected? Which actions open or close
whole parts of the graph? Where do safe and unsafe states lie relative to one
another? Which subgraphs correspond to progress, stuttering, obstruction, or
recovery?

A graph answers such questions in a form richer than a pass/fail verdict.
Checking tells us whether a property holds, but the graph reveals why the system
has the shape it does. This distinction is practical rather than rhetorical.
Once the graph exists, it can support explanation, debugging, design
comparison, teaching, and publication.

For that reason, \texttt{dstr} is best understood not simply as a specification
language or a checker, but as a graph generator for behavioral systems.

\section{Graph Query as a Second Stage}

This perspective becomes much more powerful when paired with a graph query
language. The role of \texttt{greggle} in this ecosystem is precisely that. If
\texttt{dstr} produces a labelled directed graph of reachable behavior, then
\texttt{greggle} provides a way to ask structured questions about paths,
patterns, and relational configurations inside that graph. The companion tools
\texttt{greggle} and \texttt{gruggle} are described in more detail in earlier
work on regular path queries and graph manipulation
\cite{Majumder2026GreggleGruggle}.

This pairing is compelling because the two layers address different questions.
\texttt{dstr} is about \emph{stating} a dynamic system. \texttt{greggle} is
about \emph{interrogating} the graph that emerges. One can imagine asking:

\begin{itemize}[leftmargin=*]
\item whether there exists a path from an initial state to a region marked by a
particular action pattern,
\item whether unsafe states are reachable only through some narrow corridor of
intermediate states,
\item whether two families of states are connected by alternating classes of
transitions, or
\item whether there are recurrent motifs in the structure of recovery or
escalation.
\end{itemize}

These are graph questions rather than front-end specification questions.
Treating the reachable system as a first-class graph allows analysis to proceed
at the right level of abstraction. The state-action language establishes the
world; the graph query language studies its internal geometry.

\section{Graph Manipulation as Presentation and Analysis}

After querying comes shaping. Large graphs are rarely useful in raw form. They
need to be sliced, filtered, annotated, highlighted, and projected for a human
reader. This is where \texttt{gruggle} fits naturally beside \texttt{dstr} and
\texttt{greggle}.

\texttt{gruggle} treats a graph as an object that can be manipulated without
losing its identity as a graph. Paths can be retained or discarded. Nodes can be
selected by structural criteria. Styles can be attached to the pieces that carry
the explanatory burden. A dense state space can therefore be turned into a
sequence of views: the safe core, the error boundary, the progress skeleton, the
interesting corridor, the counterexample trace, the neighborhood of a deadlock.

In this combined workflow, \texttt{dstr} supplies the behavioral universe,
\texttt{greggle} finds the meaningful substructure, and \texttt{gruggle}
produces a graph shaped for thought and communication. This is more than tool
interoperability. It is a coherent methodology:

\begin{enumerate}[leftmargin=*]
\item describe a system in a compact state-action language,
\item enumerate its reachable behavior,
\item interrogate the resulting graph with structured path queries, and
\item manipulate the graph into explanatory views.
\end{enumerate}

\section{Design Commitments}

The ambition of \texttt{dstr} is not to be maximal. It is to make a particular
design space tractable and productive. The language is intentionally small,
finite-state oriented, and explicit. Those constraints are not admissions of
defeat. They are what make the whole pipeline concrete.

Because the models are finite and enumerable, the resulting graphs are not
hypothetical semantic objects. They are actual artifacts that can be exported,
visualized, queried, versioned, and compared. Because the DSL is succinct, the
cost of expressing another system, another invariant, or another variant remains
low. Because the intermediate form is simple, the tooling around it can remain
transparent.

This is a useful stance for exploratory formal work. Rather than beginning with
a grand universal formalism, one begins with an effective language for stating
systems that matter, and then insists that the semantics become explicit enough
to support downstream graph reasoning.

\section{Use Cases}

\texttt{dstr} is especially well suited to situations where one wants both
formal discipline and graph-level understanding:

\begin{itemize}[leftmargin=*]
\item coordination protocols whose failure modes are easier to see as paths and
regions than as isolated traces,
\item controller logic where the state space is finite but the interaction among
modes is subtle,
\item workflow and policy systems where reachability, blockage, and recoverable
states matter,
\item teaching settings where students benefit from seeing how behavioral rules
compile into a concrete transition graph, and
\item research settings where the state graph is not the endpoint but the input
to a richer graph-analytic pipeline.
\end{itemize}

The last case is especially important. Once behavioral models are routinely
turned into graphs, new forms of analysis become natural: graph differencing
between design revisions, graph querying for structural motifs, graph projection
for role-specific views, and graph rewriting for what-if transformations.

\section{A Graph-Centered View of Formal Modeling}

There is a broader philosophical claim underneath \texttt{dstr}. Formal modeling
is often imagined as a narrow act of verification: write a formal description,
ask whether a theorem holds, and stop there. But a formal model can be more than
an instrument for yes-or-no answers. It can be a generative source of structured
objects that participate in a larger computational ecosystem.

In \texttt{dstr}, the formal model generates a graph. That graph can then be
handled by general graph tools, not merely by the checker that produced it. This
creates a bridge between behavioral modeling and graph computation. The benefit
is conceptual as much as practical. It invites us to regard state spaces as
objects of ongoing inquiry rather than hidden machinery inside verification.

\section{Conclusion}

\texttt{dstr} proposes a simple but powerful workflow. Start with a concise
s-expression language for describing state, action, and property. Normalize that
description into a stable machine form. Enumerate the reachable behavior
explicitly. Treat the resulting state graph as a first-class artifact. Then use
query and manipulation tools such as \texttt{greggle} and \texttt{gruggle} to
study and present that graph with far greater precision than a checker report
alone can offer.

Seen this way, \texttt{dstr} is not merely a compact formalism for finite
systems. It is an entry point into a larger graph-centered style of reasoning
about behavior. That style is attractive because it keeps three virtues in view
at once: the expressive clarity of the surface DSL, the semantic firmness of
explicit state exploration, and the analytic freedom of downstream graph query
and graph manipulation.

\bibliographystyle{plain}
\bibliography{dstr_paper}

\end{document}

